\section{联合熵、条件熵与链式法则}
\label{sec:07-Information-Theory/02-Joint-Information/01}

想同时描述两个人的行踪：A 今天去了哪，B 今天去了哪。若两人毫无关系，可以各自独立描述，总信息量就是两人的熵之和。但若 A 和 B 是同事、同一个办公室，描述了 A 的位置后，B 的位置已经八九不离十——再单独用同样的代价去描述 B，显然浪费。

两种信息的加和规则，不是简单的加法。

\subsection{联合熵：把一对变量当整体}

最直接的办法是将 $(X,Y)$ 视为一个联合随机变量：

\[
H(X,Y)=-\sum_{x,y}p(x,y)\log p(x,y).
\]

把这个定义往两个极限方向推一推，直觉就清楚了。

若 $X$ 和 $Y$ 独立，$p(x,y)=p(x)p(y)$，$\log$ 拆开，求和也拆开：

\[
H(X,Y)=H(X)+H(Y).
\]

互不相关则信息是加性的——这符合预期。

若 $Y=X$——两份记录其实是同一件事：

\[
H(X,Y)=H(X)=H(Y).
\]

描述了 $X$ 之后，$Y$ 没有带来任何新信息，联合熵不增加。这条极值不是定义给出的，是代入 $p(x,y)=p(x)\delta_{y,x}$ 算出来的。

一般情况则落在这两个边界之间：

\[
\max\{H(X),H(Y)\}\le H(X,Y)\le H(X)+H(Y).
\]

下界表示：联合描述的信息量，至少不比单独描述任何一个少（你毕竟同时描述了 $X$ 和 $Y$）。上界表示：最多也就是各说各的，独立性是信息相加的上限。两个界的理由来自本节后面要展开的链式法则和条件化性质。

\subsection{条件熵：已知 Y 后 X 还剩多少不确定}

给定 $Y=y$ 这一具体观察后，$X$ 的条件分布为 $p(x\mid y)$。沿着这个条件分布再写一次熵：

\[
H(X\mid Y=y)=-\sum_x p(x\mid y)\log p(x\mid y).
\]

再按所有可能的 $y$ 取加权平均：

\[
H(X\mid Y)=\sum_y p(y)\,H(X\mid Y=y).
\]

这就是条件熵——在知道 $Y$ 的结果之后，平均而言对 $X$ 还剩多少不确定。

这里有一个反直觉但关键的细节。不是每个 $y$ 都让你更确定 $X$。某次具体的观察可能反而让情况更混乱：

\[
H(X\mid Y=y)>H(X)
\]

完全可能。一个医学检测结果也许排除了一些可能性，但又打开了另一些罕见但严重的诊断——观察前后，$X$ 的条件熵可能比无条件熵还大。只有在所有可能的 $y$ 上取平均，才有

\[
H(X\mid Y)\le H(X).
\]

``条件化平均减少熵''是一个平均陈述，不是对每个观察值逐点担保。不要把平均性质和逐点性质搞混——这是初学条件熵时最容易踩的坑。

\subsection{链式法则：把联合描述拆成顺序步骤}

从 $p(x,y)=p(y)p(x\mid y)$ 出发。两边取 $-\log$，再对联合分布求期望。$-\log p(y)$ 的期望是 $H(Y)$；$-\log p(x\mid y)$ 对联合分布的期望正是 $H(X\mid Y)$（写出双重求和再按定义归并即可验证）。因此

\[
H(X,Y)=H(Y)+H(X\mid Y).
\]

交换角色也有对称版本：

\[
H(X,Y)=H(X)+H(Y\mid X).
\]

推广到多个变量：

\[
H(X_1,\ldots,X_n)=\sum_{i=1}^{n} H(X_i\mid X_1,\ldots,X_{i-1}).
\]

这恰好是概率链式法则在 $-\log$ 期望下的影子：$p(x_1,\ldots,x_n)=\prod_i p(x_i\mid x_1,\ldots,x_{i-1})$ 取负对数再平均。

回到日常的例子。设 $Y$ 记录当天是否下雨，$X$ 记录你是否带伞。单独描述``是否带伞''的熵可能不小——你不总是按天气做决定。但若先知道天气，再补充伞的信息，所需的额外描述短得多。要完整描述两件事，可以按顺序：先描述天气（$H(Y)$），再在天气已知后补充带伞选择（$H(X\mid Y)$）。两种方式——先天气再伞、还是先伞再天气——最终联合信息量相同。等号不要求伞完全由天气决定，也不声称天气是唯一原因；它只是用信息量语言把联合不确定性按描述顺序拆开。

\subsection{链式法则碰上一个硬例子}

视频预测编码是链式法则的工程肉身。单独一帧像素的熵往往很高；但给定前几帧和运动补偿预测后，残差的熵常常小得多。编码器发送少量 I 帧（参考帧），然后发送相对于预测的增量信息。这就是

\[
H(\text{视频序列})=H(\text{参考帧})+H(\text{残差}\mid\text{参考帧和运动预测}).
\]

当镜头突然切换、快速运动或传感器噪声涌入时，残差变大，码率随之增加。链式法则本身不担保随便什么预测器都能压缩——错模型消除不掉信息，只把未解释的结构遗留在残差和模型开销里。链式法则写的是信息关系，不是任何特定算法的性能保证。

\subsection{用手算一个例子：带噪复制}

令 $X\sim\mathrm{Bernoulli}(1/2)$，独立的噪声 $N\sim\mathrm{Bernoulli}(\varepsilon)$，观测 $Y=X\oplus N$（模二加）。

$X$ 公平，$Y$ 也公平（加独立噪声后 0 与 1 仍然各半）：

\[
H(X)=H(Y)=1\ \text{bit}.
\]

给定 $X$ 后，$Y$ 的不确定性只剩噪声：

\[
H(Y\mid X)=h_2(\varepsilon),
\]

$h_2$ 是上一章的二元熵函数。因此

\[
H(X,Y)=H(X)+H(Y\mid X)=1+h_2(\varepsilon).
\]

当 $\varepsilon=0$：$Y=X$，联合熵为 1 bit——两个变量完全重复；当 $\varepsilon=1/2$：$Y$ 与 $X$ 独立，联合熵为 2 bit——两份独立的一 bit 信息。

\subsection{连续版本的条件微分熵}

连续情形照搬定义：

\[
h(X,Y)=-\iint f(x,y)\log f(x,y)\,dx\,dy,
\]

并在合理条件下设

\[
h(X\mid Y)=h(X,Y)-h(Y).
\]

条件微分熵可以为负——这与微分熵可以为负是一脉相承的。更稳健的做法是用 KL 散度直接定义互信息：

\[
I(X;Y)=D(P_{XY}\|P_XP_Y)\ge 0,
\]

当相关微分熵都定义良好且差值不出现 $\infty-\infty$ 时，再写成

\[
I(X;Y)=h(X)-h(X\mid Y).
\]

连续情形中，应更多依赖互信息和 KL 散度，不要把一个单独的微分熵值当作绝对 bit 数来解读，也不要仅凭形式上的熵差就断言非负。
